% =========================================================================
% Capítulo 3: Cinética de sistemas de partículas
% Tema 01: Sistemas de partículas
% Problema 12: Péndulo acoplado a soporte con resorte horizontal
% =========================================================================
\subsubsection{Péndulo acoplado a soporte con resorte horizontal}

Se considera un sistema compuesto por un carro de masa $M$ que se desplaza sin fricción sobre una guía horizontal. El carro está unido a una pared fija mediante un resorte lineal de constante elástica $k$ y longitud natural $l_0$. Desde el centro del carro, en el punto $O$, pende una varilla rígida de longitud $L$ y masa despreciable, en cuyo extremo se encuentra una partícula de masa $m$. El movimiento del carro se describe mediante la coordenada horizontal $x(t)$, mientras que la orientación del péndulo se define por el ángulo $\theta(t)$ medido respecto a la vertical, como se muestra en la \textbf{Figura \ref{fig:pendulo_resorte_movil}}. Se asume que la gravedad $g$ actúa en dirección vertical hacia abajo.

\begin{enumerate}
    \item Defina los marcos de referencia y sistemas de coordenadas adecuados para describir la configuración del sistema.
    \item Determine y analice la cinemática de la partícula $m$ y de la masa $M$ en términos de las coordenadas generalizadas $\{x, \theta\}$.
    \item Determine y analice la cinética del sistema para obtener las ecuaciones diferenciales de movimiento.
\end{enumerate}

\begin{figure}[!ht]
    \centering
    \begin{tikzpicture}[>=Latex, scale=1.3]
        % --- CONFIGURACIÓN DE PARÁMETROS ---
        \def\xCar{1.5}      % Posición x del carro
        \def\ang{30}       % Ángulo theta
        \def\L{2.5}        % Longitud del péndulo
        \def\wCar{1.0}     % Ancho del carro
        \def\hCar{0.6}     % Alto del carro

        % --- 1. ENTORNO (Suelo y Pared) ---
        \fill[pattern=north east lines] (-1.5, -0.2) rectangle (4, 0);
        \draw[thick] (-1.5, 0) -- (4, 0);
        \fill[pattern=north east lines] (-1.7, 0) rectangle (-1.5, 1.5);
        \draw[thick] (-1.5, 0) -- (-1.5, 1.5);

        % --- 2. RESORTE ---
        \draw[decoration={aspect=0.3, segment length=4pt, amplitude=3pt, coil}, decorate] 
            (-1.5, \hCar/2) -- (\xCar-\wCar/2, \hCar/2);
        \node[above] at ({(-1.5 + \xCar-\wCar/2)/2}, \hCar/2 + 0.1) {$k$};

        % --- 3. CARRO (Masa M) ---
        \draw[fill=gray!10, thick] (\xCar-\wCar/2, 0.1) rectangle (\xCar+\wCar/2, \hCar+0.1);
        \draw[fill=black] (\xCar-0.3, 0.1) circle (0.1);
        \draw[fill=black] (\xCar+0.3, 0.1) circle (0.1);
        \node[left] at (\xCar, \hCar/2+0.1) {$M$};
        \coordinate (O) at (\xCar, \hCar/2+0.1);
        \filldraw (O) circle (1.5pt) node[above right] {$O$};

        % --- 4. PÉNDULO (Masa m) ---
        \begin{scope}[shift={(O)}, rotate=\ang]
            \draw[ultra thick] (0,0) -- (0, -\L);
            \draw[line width=2.5pt,line cap= round] (0,0) -- (0, -\L);
            \draw[line width=1.5pt, gray,line cap= round] (0,0) -- (0, -\L);
            \shadedraw[ball color=gray!30] (0, -\L) circle (0.25) node[] {$m$};
            \node[right] at (0,-\L/2) {$L$};
        \end{scope}

        % --- 5. REFERENCIAS Y ÁNGULOS ---
        % Eje vertical de referencia
        \draw[dashed, gray] (O) -- (\xCar, -\L-0.5);
        % Ángulo theta
        \draw[->] (\xCar, {\hCar/2+0.1-0.8}) arc (270:{270+\ang}:0.8);
        \node at (\xCar+0.3, {\hCar/2+0.1-1.0}) {$\theta$};
        
        % Cota x(t)
        \draw[dashed, thin] (0, 1.2) -- (0, 0);
        \draw[|->, thick] (0, 1.0) -- (\xCar, 1.0) node[midway, above] {$x(t)$};

        % Vector Gravedad
        \draw[->, thick] (3.5, 1.5) -- (3.5, 0.8) node[below] {$\mathbf{g}$};

    \end{tikzpicture}
    \caption{Diagrama de un péndulo de masa $m$ acoplado a un soporte móvil de masa $M$ con restricción elástica.}
    \label{fig:pendulo_resorte_movil}
\end{figure}